\section{Theoretical results}
\label{sec:theory}

Three partial results follow; their precise statements, proofs and scope are in Section~S2 of
the Supplement, and none covers the passage from one stage of ESF to the next, the weighted
vote, or the search for the candidate line. Throughout, $O\subseteq\{1,\dots,n\}$ is a fixed set of
contaminated units whose values are arbitrary, and a subsample is \emph{clean} if it contains
no unit of $O$.

\paragraph{One stage of the sequential draw.}
\label{sec:stage}
Let $E_\ell$ be the set of units not flagged after stage $\ell-1$, from which stage $\ell$
draws, and $\varepsilon_\ell=|E_\ell\cap O|/|E_\ell|$.

\begin{proposition}[informal statement; the precise version is in Supplement~S2.0]
\label{prop:clean}
Conditionally on the data and on everything drawn before stage $\ell$, the subsamples of stage
$\ell$ are independent and uniform on the admissible subsamples of size $m$ of $E_\ell$, each
replicate is clean with probability $\pi_\ell=p_\ell a_\ell$, where
$p_\ell\le(1-\varepsilon_\ell)^{m}$ is the probability that a uniform subsample of $E_\ell$ is
clean and $a_\ell$ corrects for admissibility, and the stage contains a clean replicate
with probability $1-(1-\pi_\ell)^{B/L}$.
\end{proposition}

The result is exact and covers one stage: a clean replicate becomes likelier as the flags
remove outliers from $E_\ell$, but nothing here says that earlier stages do so. With $m=8$ the
first stage contains a clean replicate with probability about $0.84$ at $\varepsilon=0.20$ and
$0.45$ at $0.30$, and with $m=12$ about $0.51$ and $0.13$, which is where ESF breaks down.

\paragraph{The final reweighting step.}
\label{sec:reweight}
The reweighting step takes a parameter $\theta$, recomputes the flags and refits each group on
its unflagged units. Let the units be independent with law $P=(1-\varepsilon)P_0+\varepsilon C$,
where under $P_0$ a unit of group $k$, drawn with probability $\pi_k>0$, has
$y=x^{\top}\theta^{0}_k+\sigma_ke$ with $e\sim\mathcal N(0,1)$ independent of $x$, and every
unit of $C$ lies beyond $c'>c$ scales from the nearest line; let $\eta$ be the probability that
a clean unit is nearer to another true line than to its own.

\begin{proposition}[informal statement; the precise version is Theorem~S2 of the Supplement]
\label{thm:step}
Apply the map at $\theta=\theta^{0}$. Under mild regularity conditions, almost surely as
$n\to\infty$, the scale of group $k$ converges to a limit $s_k$ with $s_k/\sigma_k=1+O(\nu_k)$,
$\nu_k$ being of the order of $\eta+\varepsilon$ divided by the mass assigned to line $k$; the
flagged fraction converges to $\varepsilon+(1-\varepsilon)\kappa$, where $\kappa$, the fraction
of clean units flagged, lies within $\eta$ of $\sum_k\pi_k2\Phi(-cs_k/\sigma_k)$; and the
refitted coefficients converge to a limit that depends on the contamination only through the
$s_k$ and lies within $O(\eta)$ of $\theta^{0}_k$ when the covariates are bounded. Without
overlap or contamination, $\theta^{0}$ is a fixed point of the map and the flagged fraction
tends to $2\Phi(-c)$.
\end{proposition}

This is a law of large numbers for one application of the map at a fixed input: it says what
the flagged fraction estimates under Gaussian errors, $\varepsilon$ plus a fraction of clean
units close to $2\Phi(-c)\approx0.012$ when the groups overlap little, and that contamination
beyond the cut-off is removed entirely. It does not cover four things: a small group under
moderate contamination, where the condition $\nu_k<\tfrac12$ of Theorem~S2 fails; the input
that ESF delivers to the map; the repeated application of the map in the code; and the flagged
fraction reported at the refitted parameter (Supplement, Section~S2.2).

\paragraph{The group-recovery step.}
\label{sec:recovery-theory}
The recovery step scores a \emph{fit} $F$, a set of lines $\theta_k$ with scales $s_k$ and
proportions $\pi_k$ and a noise weight $\pi_0>0$ summing to one, by
$\ell_n(F)=\sum_{i\le n}\log f_F(u_i)$, where
$f_F(u)=\sum_{k}(\pi_k/s_k)\,\varphi\{(y-x^{\top}\theta_k)/s_k\}+\pi_0/R$. For given
parameters, $\ell_n(F)/n$ converges almost surely to $L(F)=\E_P\log f_F(u)$; the proposition
values, in $L$, the two operations by which a candidate differs from the original fit.

\begin{proposition}[informal statement; the precise version is in Supplement~S2.3]
\label{prop:recover-value}
(i) Let $G$ be an event of probability $\pi_\ast$ on which $y=x^{\top}\theta_\ast+\sigma_\ast e$
with $e\sim\mathcal N(0,1)$ independent of $x$, let $F$ have noise weight $\pi_0>\pi_\ast$, and
let $F'$ be $F$ with the line $(\theta_\ast,\sigma_\ast,\pi_\ast)$ added and the noise weight
lowered to $\pi_0-\pi_\ast$. Then
\begin{equation}
L(F')-L(F)\;\ge\;\pi_\ast\Bigl\{\log\frac{\pi_\ast R}{\pi_0}-\log\sigma_\ast
-\tfrac12\log(2\pi e)-\xi\Bigr\}-\nu\log\frac{\pi_0}{\pi_0-\pi_\ast},
\label{eq:recover-gain}
\end{equation}
where $\xi\ge0$ is small when $G$ lies a few scales from every line of $F$, and
$\nu\le1-\pi_\ast$ is the noise mass that $F$ attributes to the units outside $G$. (ii) If $F$
contains two lines through one group of probability $\pi_S$ and scale $\sigma_S$, at intercepts
$\delta$ apart, and $F''$ replaces them by the line through the group, then
$|L(F'')-L(F)|\le(\varphi(1)a^{2}+\tfrac14a^{4})\pi_S$ with $a=\delta/(2\sigma_S)$.
\end{proposition}

The leading term of \eqref{eq:recover-gain} is the value of a group: per unit, the log of
the ratio of its peak density $\pi_\ast/\sigma_\ast$ to the noise density $\pi_0/R$, less the
entropy $\tfrac12\log(2\pi e)\approx1.42$ of a unit normal error. With $R/\sigma_\ast=40$,
$\pi_\ast=0.1$, $\pi_0=0.2$ and $\xi\approx0$, the value is $1.58$ nats per unit of the group,
that is $0.158n$ in all. With these values the loss term is $\nu\log2$: for a fit whose noise
mass off $G$ is about $0.1$, the weight $\pi_0=0.2$ less the group's $0.1$, it is about $0.069n$
and the gain at least $0.09n$ against a margin $\tfrac{d+1}{2}\log n$ of $8.6$ for $d=2$,
$n=300$. The bound $\nu\le1-\pi_\ast$ alone gives a loss of up to $0.62n$, which exceeds the
value, so the inequality is informative only for such a fit. Part (ii), an idealisation with two parallel lines at intercepts $\delta$ apart, says that
removing one of two lines through one group is nearly free, so a candidate that arises by the two operations is accepted
for all large $n$ when the right side of \eqref{eq:recover-gain} exceeds $\varphi(1)a^{2}\pi_S$
and its scales are admissible. The result is for fixed parameters, which the step reads off
the data, and it does not cover the reweighting refit of the candidate. The peak test, analysed
in Supplement~S2.3, fails a band of outliers that fills its window and passes one of half-width
below about $3.4$ scales.

